\clearpage
\VerdierPhysicalPageAnchor{0044}%
\VerdierSemanticAnchor{ch1-1-4}
\subsection{群の変更}

\VerdierSemanticAnchor{ch1-1-4-1}
\noindent\textbf{1.4.1.}\quad
$f\colon H\to G$ を群準同型とする。各 $G$-圏 $(\mathcal C,G,T,c)$ に対し、
$f$ を介して $H$ を $\mathcal C$ に作用させる。すなわち

\VerdierSemanticAnchor{ch1-1-4-1-eq1}
\begin{equation}
  \begin{aligned}
    T_f(h)&=T\bigl(f(h)\bigr),
      & h&\in H,\\
    c_f(h,h')&=c\bigl(f(h),f(h')\bigr),
      & h,h'&\in H
  \end{aligned}
  \tag{1.4.1.1}
\end{equation}
とおくことにより、$H$-圏 $(\mathcal C,H,T_f,c_f)$ を対応させる。

\VerdierSemanticAnchor{ch1-1-4-1-gfunctor}
同様に、$G$-関手
\[
  (F,m)\colon(\mathcal C,G,T,c)
  \longrightarrow(\mathcal C',G,T',c')
\]
が与えられたとき、

\VerdierSemanticAnchor{ch1-1-4-1-eq2}
\begin{equation}
  m_f(h)=m\bigl(f(h)\bigr)
  \tag{1.4.1.2}
\end{equation}
とおいて、対応する $H$-圏の間の $H$-関手 $(F,m_f)$ を対応させる。

\VerdierSemanticAnchor{ch1-1-4-1-morphisms}
最後に、$G$ の作用と両立する $G$-関手の射は $H$ の作用とも両立する。

\VerdierSemanticAnchor{ch1-1-4-1-eq3}
以上は明らかに一つの関手を定める。これを\emph{群の変更関手}という：
\begin{equation}
  f^*\colon G\text{-}\operatorname{Cat}
  \longrightarrow H\text{-}\operatorname{Cat}.
  \tag{1.4.1.3}
\end{equation}

\VerdierSemanticAnchor{ch1-1-4-1-eq4}
より正確には、$\mathcal C,\mathcal C'$ を二つの $G$-圏とすると、関手 $f^*$
は写像
\begin{equation}
  \begin{aligned}
    f^*(\mathcal C,\mathcal C')\colon\quad
      &\operatorname{Hom}_{G\text{-}\operatorname{Cat}}
        (\mathcal C,\mathcal C')\\
      &\longrightarrow
        \operatorname{Hom}_{H\text{-}\operatorname{Cat}}
        \bigl(f^*(\mathcal C),f^*(\mathcal C')\bigr)
  \end{aligned}
  \tag{1.4.1.4}
\end{equation}
を定める。

\VerdierSemanticAnchor{ch1-1-4-1-homcategories}
しかし集合\allowbreak
$\operatorname{Hom}_{G\text{-}\operatorname{Cat}}(\mathcal C,\mathcal C')$\allowbreak
と\allowbreak
$\operatorname{Hom}_{H\text{-}\operatorname{Cat}}
\bigl(f^*(\mathcal C),f^*(\mathcal C')\bigr)$\allowbreak
は、それぞれ圏\hfill\break
$\mathcal H\!om_{G\text{-}\operatorname{Cat}}(\mathcal C,\mathcal C')$\allowbreak
と\allowbreak
$\mathcal H\!om_{H\text{-}\operatorname{Cat}}
\bigl(f^*(\mathcal C),f^*(\mathcal C')\bigr)$\allowbreak
の対象集合であり、その射は群作用と両立する関手の射である。また写像
$f^*(\mathcal C,\mathcal C')$ は自然に関手

\VerdierSemanticAnchor{ch1-1-4-1-eq5}
\begin{equation}
  \begin{aligned}
    f^*(\mathcal C,\mathcal C')\colon\quad
      &\mathcal H\!om_{G\text{-}\operatorname{Cat}}
        (\mathcal C,\mathcal C')\\
      &\longrightarrow
        \mathcal H\!om_{H\text{-}\operatorname{Cat}}
        \bigl(f^*(\mathcal C),f^*(\mathcal C')\bigr)
  \end{aligned}
  \tag{1.4.1.5}
\end{equation}
へ拡張される。

\VerdierSemanticAnchor{ch1-1-4-1-faithful}
この関手は忠実である。
